\chapter{Conclusiones}
\label{chapter:8-Conclusiones}

Se logró desarrollar satisfactoriamente un chatbot educativo basado en el contenido de OpenFing. Mediante nuestra prueba de concepto logramos mostrar que el alto volumen de datos, el gran cómputo requerido y el tiempo que demanda el preprocesamiento de las clases no es un impedimento. A pesar de las dificultades mencionadas, la solución es viable y funciona como es esperada.

Los resultados concretos obtenidos respaldan esta conclusión. En la evaluación de transcripción automática, el modelo Whisper \textit{medium} con especificación explícita del idioma español alcanzó una WER promedio inferior a 0.05 en todos los conjuntos de prueba evaluados, equivalente a menos de cinco errores por cada cien palabras. El pipeline de transcripción fue aplicado al corpus completo de OpenFing, que comprende 239 cursos y 4.339 clases, utilizando \textit{faster-whisper} con CTranslate2, lo que redujo los tiempos de procesamiento entre dos y cuatro veces respecto a la implementación original de Whisper. La generación de respuestas se implementó con el modelo \textit{meta-llama/Meta-Llama-3.1-8B-Instruct} (Llama~3.1~8B) accedido mediante la API de inferencia de HuggingFace, que demostró ser adecuado para el dominio académico en español sin requerir infraestructura de GPU propia.

Para el estudiante, la solución brinda respuestas mayoritariamente correctas con relación a las clases y cursos en los que se basa, en tiempo y forma razonable. Puede resultar muy útil a la hora de estudiar, ya que facilita el aprendizaje, agiliza la asimilación de conceptos y apoya el estudio autónomo, consiguiendo respuestas inmediatas sobre consultas específicas y complementando los recursos actuales.

Para la facultad, esta prueba de concepto representa una oportunidad para escalar el acceso al conocimiento. La implementación logró consolidar material educativo disperso en los múltiples videos en una interfaz conversacional única. Los recursos computacionales y de procesamiento necesarios para mantener el proyecto ya están disponibles en la infraestructura de la facultad, de modo que su operación continua no debería representar una carga significativa.

Respecto a los objetivos específicos planteados en la \autoref{section:1-Objetivos del Proyecto}, el proyecto los abordó en su totalidad aunque con distinta profundidad. El primer objetivo, diseñar e implementar un pipeline automatizado de extracción, descarga, transcripción e indexación, fue cumplido con un sistema modular de cinco componentes orquestados mediante Prefect. El segundo objetivo, evaluar herramientas \acrshort{asr} para el español académico rioplatense, fue satisfecho con la evaluación cuantitativa y cualitativa de los modelos Whisper \textit{base}, \textit{small} y \textit{medium}, concluyendo en la selección de \textit{medium} con especificación de idioma. El tercer objetivo, implementar un chatbot \acrshort{rag} con referencia al instante exacto del video, fue cumplido: el sistema recupera fragmentos de transcripción e incluye metadatos de temporización en cada respuesta. El cuarto objetivo, evaluar la calidad del pipeline \acrshort{rag} mediante experimentos de pregunta-respuesta, fue abordado parcialmente: la evaluación de transcripción quedó completa y documentada, mientras que la evaluación del componente conversacional sobre el corpus completo constituye la principal área pendiente. El quinto objetivo, desarrollar el sistema como software de código abierto con arquitectura extensible, fue logrado con una arquitectura de microservicios que permite incorporar nuevas bases de conocimiento sin rediseño estructural.

Si bien el prototipo actual cumple su función demostrativa, existen áreas claras de mejora: refinamiento de la precisión de respuestas, implementación de métricas de uso, integración con sistemas académicos existentes, y expansión del corpus con material de diversas asignaturas. Estas mejoras son técnicamente factibles y mantienen la filosofía de eficiencia de recursos del proyecto original.

En términos de autocrítica, el alcance del procesamiento resultó más amplio de lo inicialmente estimado. El preprocesamiento del corpus completo de OpenFing demandó un tiempo de cómputo significativo que condicionó los ciclos de experimentación con distintas configuraciones. Una estimación más precisa de los costos computacionales en la etapa de planificación habría permitido distribuir mejor el tiempo disponible. Por otra parte, la evaluación del componente de pregunta-respuesta no alcanzó la profundidad prevista, ya que el tiempo disponible priorizó la construcción y estabilización del pipeline completo sobre la fase experimental.

El proyecto valida que las barreras para implementar herramientas educativas basadas en \acrshort{ia} son principalmente metodológicas, no económicas, abriendo posibilidades para democratizar el acceso a asistentes pedagógicos inteligentes en contextos académicos.
